\section{4.1 High-Level Architecture}
The system architecture has four pipelines:
\begin{enumerate}
  \item Calibration pipeline (intrinsics + extrinsics).
  \item Acquisition pipeline (synchronized 4-camera capture).
  \item Inference and triangulation pipeline (ball + joints).
  \item Evaluation and visualization pipeline.
\end{enumerate}

All pipelines operate in a shared world frame defined by \texttt{Dimensions.txt} and extrinsics JSON.

\section{4.2 Camera Projection Model}
For each camera, the projection model is:

\[ s\mathbf{u} = \mathbf{K}[\mathbf{R}|\mathbf{t}]\mathbf{X}_w \]

where:
- \(\mathbf{X}_w\) is a 3D world point,
- \(\mathbf{R}, \mathbf{t}\) are extrinsics,
- \(\mathbf{K}\) is intrinsics matrix,
- \(\mathbf{u}\) is image point.

Distortion is handled by OpenCV distortion coefficients and undistortion routines. In `process_4cam_to_3d.py`, points are undistorted before triangulation to operate in normalized coordinates.

## 4.3 Triangulation Formulation
Given normalized observations \((x_i, y_i)\) and projection matrix \(P_i\), DLT constructs a linear system:

\[
A\mathbf{X}=0,
\]

with two equations per camera:
\[
x_iP_{i,3} - P_{i,1}=0,
\]
\[
y_iP_{i,3} - P_{i,2}=0.
\]

The solution is obtained by SVD, using the right singular vector corresponding to the smallest singular value, then homogeneous normalization.

\section{4.4 Robust Reprojection-Based Rejection}
Naive triangulation from all cameras is vulnerable to one bad observation. The robust logic in this thesis triangulates then reprojects to each camera and computes pixel error. If worst error exceeds threshold, that camera is dropped and triangulation repeats until error is acceptable or minimum camera count is violated.

This iterative rejection is central for practical robustness, especially during dynamic ball tests.

\section{4.5 Temporal Filtering for Ball Trajectory}
Optional post-triangulation filters include:
\begin{itemize}
  \item speed gate: reject physically implausible jumps,
  \item EMA smoothing: reduce jitter while preserving trajectory trend.
\end{itemize}

These filters are tuned conservatively to avoid over-smoothing true dynamics. The thesis explicitly compares raw and filtered behavior to prevent hidden bias.

\section{4.6 2D Pose to 3D Joint Triangulation}
For each frame and each camera:
\begin{enumerate}
  \item run person detection + keypoint model,
  \item select target person with identity-consistency heuristic,
  \item for each joint, gather confident 2D observations from multiple cameras,
  \item triangulate joint 3D with minimum-camera requirement.
\end{enumerate}

Per-joint triangulation ensures one bad keypoint does not invalidate all joints.

\section{4.7 Coordinate Frame Definition}
The arena frame is expressed in millimeters with origin at north-east floor corner. Wall definitions and AprilTag corner coordinates are listed in \texttt{garage\_lab\_combined/cal/extrinsics/Dimensions.txt}. Camera positions measured in the same frame provide physically interpretable outputs.

Consistency in units is critical. Earlier mixed-unit drafts (cm/mm) were removed and unified to mm across GT, processing outputs, and reporting.

\section{4.8 Calibration Methodology}
\subsection{4.8.1 Intrinsics}
Intrinsics are solved per camera using ChArUco images captured at runtime resolution (1280x720). Frames with low corner quality are excluded via minimum-corner thresholds.

\subsection{4.8.2 Extrinsics}
Extrinsics are estimated with AprilTag detections against known tag corner coordinates. Robust options include:
\begin{itemize}
  \item point error filtering,
  \item sigma-based outlier rejection,
  \item minimum points threshold,
  \item camera-specific include-tags map.
\end{itemize}

When one camera physically moves, partial recalibration is performed for that camera and merged into final extrinsics.

\section{4.9 Synchronization Strategy}
The pipeline supports both software-based offset estimation and manual flash-based alignment. In deployment, manual flash alignment was preferred for interpretability. Aligned clips are trimmed to common start and length before 3D processing.

\section{4.10 Measurement Model for GT Evaluation}
For each trial:
\begin{itemize}
  \item estimated point \(\hat{p}\) is compared to ground truth \(p\),
  \item error vector \(e=\hat{p}-p\),
  \item norm \(|e|=\sqrt{e\_x^2+e\_y^2+e\_z^2}\).
\end{itemize}

Aggregate metrics include mean, median, RMSE, P90/P95, max, axis-wise bias, and static precision (window standard deviation).

\section{4.11 Correction Model}
Axis-wise linear correction is modeled as:

\[
x_c = a_x x + b_x,
\]
\[
y_c = a_y y + b_y,
\]
\[
z_c = a_z z + b_z.
\]

The model is fit on static GT data. It removes systematic bias but does not fix local nonlinear distortions caused by calibration errors or visibility limitations.

\section{4.12 Design Rationale}
The architecture prioritizes:
\begin{itemize}
  \item transparency over black-box fusion,
  \item metric consistency over visual-only quality,
  \item staged safety over aggressive automation.
\end{itemize}

This is aligned with the thesis goal: produce a defensible research foundation for later launcher control.


\section{4.13 Sensitivity Analysis: Why Small Geometry Errors Cause Large Targeting Errors}
A practical sensitivity observation from this project is that a visually small camera rotation can generate large target displacement in reconstructed space, especially near room edges and for points seen by only two cameras. This is consistent with triangulation geometry: when ray intersection angles are shallow, small angular errors amplify depth uncertainty.

For launcher-control relevance, this means calibration quality cannot be summarized only by average reprojection error. Spatially varying uncertainty must be considered. Points near weak-overlap regions may have significantly higher targeting risk than central points.

\section{4.14 Observability and Camera Contribution}
The system logs the number of cameras used per triangulated point. This statistic is essential because:
\begin{itemize}
  \item 2-camera triangulation is geometrically valid but less robust,
  \item 3-4 camera support allows outlier rejection and better depth stability.
\end{itemize}

Therefore, observability itself is treated as a quality feature. Future controllers should incorporate this value into confidence gating.

\section{4.15 Human Keypoint Semantics and Measurement Ambiguity}
Unlike a rigid marker, human joints in 2D keypoint models represent anatomical landmarks estimated from appearance. Shoulder location may shift with clothing, arm rotation, and partial occlusion. Knee estimation changes with leg orientation and camera angle.

Hence joint-touch GT error includes two components:
\begin{enumerate}
  \item geometric reconstruction error,
  \item semantic labeling uncertainty of the joint in 2D model space.
\end{enumerate}

This reinforces the recommendation to use rigid calibration targets for final controller calibration and treat human-joint metrics as performance on realistic but noisy semantic targets.

\section{4.16 Why the Methodology Is Appropriate for Research-Based Assessment}
The selected methodology aligns with MSc research criteria:
\begin{itemize}
  \item clearly justified methods,
  \item explicit assumptions and constraints,
  \item objective test protocols,
  \item critical evaluation of results,
  \item clear separation between completed and pending system blocks.
\end{itemize}

This creates a defensible manuscript for a research university context and avoids overstatement.



\section{4.17 Additional Mathematical Notes for Controller Interface}
When the perception stack is connected to launcher control, the following uncertainties should be propagated:
\begin{itemize}
  \item triangulation covariance as function of view geometry,
  \item temporal uncertainty from frame latency,
  \item calibration uncertainty from extrinsics drift.
\end{itemize}

A simple controller-ready uncertainty score can be computed from normalized components:
\[
U = w_1\cdot	ext{reproj\_norm} + w_2\cdot	ext{cams\_deficit} + w_3\cdot	ext{temporal\_instability},
\]
where lower values indicate safer actuation conditions. This score can be used to decide no-fire states.

\section{4.18 Latency and Time Alignment Considerations}
For live operation, synchronization is not only frame-index alignment. End-to-end latency includes:
\begin{itemize}
  \item camera capture delay,
  \item decode delay,
  \item model inference time,
  \item triangulation and visualization time,
  \item command dispatch delay.
\end{itemize}

If total latency is high relative to target motion, aiming may use outdated coordinates. Therefore, future live systems should include timestamp-based synchronization and predictive extrapolation for moving targets.

\section{4.19 Prediction for Moving Targets}
For dynamic body-part targeting, point estimates can be enhanced with short-horizon prediction:
\[
\hat{p}(t+\Delta t)=p(t)+v(t)\Delta t,
\]
with velocity estimated from filtered finite differences. More advanced filters (e.g., Kalman) can improve noise handling, but must be tuned per joint behavior.

Prediction should be confidence-gated: if uncertainty is high, actuation should be suppressed.

\section{4.20 Calibration Drift Monitoring Model}
A lightweight drift monitor can run online by projecting stable tags and comparing expected vs observed corners. If median error exceeds threshold over a sustained window, system raises \texttt{calibration\_stale} flag and blocks autonomous firing.

This approach converts calibration from offline task into monitored runtime health signal.

\section{4.21 Model Selection Rationale Under Resource Constraints}
The project selected practical detector and pose models that run on available hardware. While larger models may improve raw accuracy, deployment constraints (4-camera throughput and response needs) required balanced models. The selection criterion was end-to-end system utility, not isolated benchmark ranking.

This reflects real engineering constraints where operational reliability outranks single-module peak accuracy.



\section{4.22 Uncertainty-Aware Triangulation Extensions (Future)}
A future extension is weighted triangulation where each view contributes according to a confidence-derived variance estimate. In matrix form, this can be expressed as minimizing weighted reprojection residuals. Combined with robust loss functions, this can reduce sensitivity to outlier observations while preserving deterministic behavior.

\section{4.23 Spatial Calibration Residual Mapping}
Instead of one global correction model, residuals can be mapped in 3D using radial basis functions or low-order polynomial surfaces. This allows location-dependent corrections where corner regions systematically deviate more than central regions.

Such nonlinear correction should only be applied after geometric calibration is stable; otherwise model drift can become hard to interpret.

\section{4.24 Time-Series Stability Metrics for Live Readiness}
For live operation, point-wise error is insufficient. Additional stability metrics are needed:
\begin{itemize}
  \item frame-to-frame acceleration outliers,
  \item confidence drop frequency,
  \item temporary track loss duration,
  \item recovery latency after occlusion.
\end{itemize}

These metrics can be monitored online and used to trigger safe mode.

\section{4.25 Bridging Perception and Control Timescales}
Perception may run at 10-15 FPS while launcher control loops can run faster. A bridging strategy is needed:
\begin{itemize}
  \item hold-last-valid-target with timeout,
  \item short-term prediction with uncertainty expansion,
  \item no-fire when target freshness exceeds threshold.
\end{itemize}

This prevents stale-target actuation.

\section{4.26 Mathematical Summary}
The mathematical backbone of this thesis combines:
\begin{itemize}
  \item calibrated projection,
  \item robust multi-view triangulation,
  \item quality gating via reprojection/camera count,
  \item optional temporal filtering,
  \item and post-hoc bias correction.
\end{itemize}

This stack is intentionally modular and interpretable, supporting both scientific analysis and future control integration.
